\clearpage
\section{Complete Secondary Behavioral Results}
\label{sec:diagnostics}
\subsection{Likelihood measurement}
\Tabref{tab:likelihood} reports every main and anchor likelihood estimate; \tabref{tab:likelihood-contrasts} reports all five separately corrected contrasts. Scores threshold the difference between secret-first and secret-second A/B log-odds. Exact ties receive half credit. We use explicit non-padding position IDs in forward scoring. Raw logits, token IDs, both orders, and full prompts remain in \artifact{results/model_outputs/choice_likelihoods.jsonl}. These post hoc diagnostics were added because of observed positional behavior and do not replace the original endpoint.
\begin{table}[htbp]
\centering\small
\begin{tabular}{@{}llll@{}}
\toprule
Study & Task & \textsc{Condition} & Accuracy [95\% CI] \\
\midrule
Anchor & Word & Free anchor & 54.2\% [41.7, 66.7] \\
Main & Plot & Baseline & 40.8\% [28.3, 53.3] \\
Main & Plot & Context & 42.5\% [27.5, 59.2] \\
Main & Plot & Outline & 40.8\% [26.7, 56.7] \\
Main & Word & Baseline & 62.8\% [52.8, 72.8] \\
Main & Word & Context & 72.8\% [63.3, 81.7] \\
Main & Word & Decoy & 61.1\% [51.1, 71.1] \\
Main & Word & Outline & 56.1\% [46.1, 66.1] \\
\bottomrule
\end{tabular}
\caption{Complete post hoc order-balanced likelihood results. Main tasks contain 90 word or 60 plot pairs per condition; the anchor has 60 pairs. Intervals resample word units or plot families. These scores are secondary and do not replace generated-choice outcomes. We do not rank below-chance results as improved privacy.}
\label{tab:likelihood}
\end{table}

\begin{table}[htbp]
\centering\small
\begin{tabular}{@{}llll@{}}
\toprule
Task & Contrast & Difference (pp) [95\% CI] & Holm $p$ \\
\midrule
Plot & Outline -- baseline & +0.0 [-19.2, +20.8] & 1.0000 \\
Plot & Outline -- context & -1.7 [-26.7, +24.2] & 1.0000 \\
Word & Outline -- baseline & -6.7 [-19.4, +6.1] & 1.0000 \\
Word & Outline -- context & -16.7 [-30.6, -1.7] & 0.1650 \\
Word & Decoy -- baseline & -1.7 [-13.9, +10.6] & 1.0000 \\
\bottomrule
\end{tabular}
\caption{All five post hoc likelihood contrasts, Holm-corrected within this separate family. None survives correction, including the outline--context word comparison.}
\label{tab:likelihood-contrasts}
\end{table}


\subsection{Plot forecasting and direction calibration}
\Tabref{tab:forecasts} gives both private-minus-no-secret contrasts and differences-in-differences. Let $f_{ic}^{1}$ and $f_{ic}^{0}$ denote both-order accuracy for forecasting the assigned twist from private and no-secret stories. We first compute $d_{ic}=f_{ic}^{1}-f_{ic}^{0}$ on complete pairs, then compare $d_{i,\mathrm{outline}}-d_{i,c}$ across common units. Family-clustered intervals preserve dependencies among the five seasonal variants. Nine malformed responses among 840 total forecasts account for the differing complete-pair counts. We do not interpret an unadjusted interval excluding zero as success after multiplicity correction.
\begin{table}[htbp]
\centering\small
\resizebox{\textwidth}{!}{%
\begin{tabular}{@{}lrll@{}}
\toprule
Forecast contrast & Pairs & Difference (pp) [95\% CI] & Holm $p$ \\
\midrule
Baseline private -- no secret & 57 & +8.8 [+2.6, +14.8] & 0.1960 \\
Context private -- no secret & 55 & +6.4 [-1.9, +14.3] & 0.7312 \\
Outline private -- no secret & 60 & -0.8 [-5.8, +5.0] & 1.0000 \\
Did: outline -- baseline & 57 & -8.8 [-16.4, +0.0] & 0.4076 \\
Did: outline -- context & 55 & -5.5 [-13.9, +3.6] & 0.7312 \\
\bottomrule
\end{tabular}
}
\caption{Complete plot-forecast contrasts. DiD subtracts the private-minus-no-secret forecast advantage of the comparison condition from that of outlines. Both candidate orders and both story-presence states must parse. Holm correction covers these five contrasts; intervals cluster by plot family. None is adjusted significant.}
\label{tab:forecasts}
\end{table}


The separate direction-calibration check uses word units 0--29 or plot families 0--3 to decide whether to invert the primary detector, then evaluates on the remainder. No primary condition selects an inversion. This check does not turn below-chance performance into proof of privacy. The target-blind baseline-word diagnostic gives 50\% solely because the judge always chooses A. Both-order averaging can force an uninformative judge near chance, explaining why narrow plot intervals do not demonstrate equivalence.

\subsection{Literal outcomes and conditional subsets}
Literal checks use case-insensitive whole-word matching. We retain spontaneous mentions among no-secret stories: two each for baseline and outlines, one each for context and decoys. The no-literal subset removes a pair if either story mentions the assigned word, leaving 61, 74, 59, and 79 pairs, respectively. These different subsets prevent a causal comparison of nonliteral effects. Secondary likelihood intervals include chance for every corresponding subset, even though point estimates range from 48.0\% to 58.5\%.

The literal-outcome Holm family contains only outline--baseline, outline--context, and decoy--baseline comparisons among secret-present stories. Their estimates and intervals appear in \secref{sec:results}; they must not be combined with primary or likelihood-family $p$-values to construct a confirmatory narrative. The systematic examples and mention-position analysis are saved in \artifact{results/error_analysis.md} and \artifact{results/literal_analysis.json}. Truncation and outline-uptake diagnostics are saved in \artifact{results/outline_and_truncation.json}.
